\documentclass[11pt]{article}
\usepackage[a4paper,margin=18mm]{geometry}
\usepackage{fontspec}
\setmainfont{Latin Modern Roman}
\usepackage{amsmath,amssymb,amsthm,amscd,graphicx,color}
\usepackage{xurl}
\usepackage[unicode,hidelinks]{hyperref}
\allowdisplaybreaks
\setlength\emergencystretch{3em}
\numberwithin{equation}{section}

\newtheorem{Theorem}{Theorem}[section]
\newtheorem{Corollary}[Theorem]{Corollary}
\newtheorem{Lemma}[Theorem]{Lemma}
\newtheorem{Proposition}[Theorem]{Proposition}
\newtheorem{Conjecture}[Theorem]{Conjecture}
 { \theoremstyle{definition}
\newtheorem{Definition}[Theorem]{Definition}
\newtheorem{Remark}[Theorem]{Remark}
\newtheorem{Fact}[Theorem]{Fact}}

\newcommand{\setatop}[2]{\genfrac{}{}{0pt}{}{#1}{#2}}

\newcommand{\vlam}{\boldsymbol{\lambda}}

\renewcommand{\Re}{\mathrm{Re}}
\renewcommand{\Im}{\mathrm{Im}}

\newcommand{\bx}{\bar{x}}
\newcommand{\bs}{\bar{s}}

\newcommand{\nonu}{\nonumber \\ \nopagebream}
\newcommand{\sign}{{\rm sign}}
\newcommand{\ee}{{\rm e}}
\newcommand{\ii}{{\rm i}}
\newcommand{\dd}{{\rm d}}

\newcommand{\R}{{\mathbb R}}
\newcommand{\C}{{\mathbb C}}
\newcommand{\Z}{{\mathbb Z}}
\newcommand{\N}{{\mathbb N}}

\newcommand{\cT}{{\mathcal T}}
\newcommand{\pdag}{^{\phantom\dag}}

\begin{document}
\begin{center}\Large For N>=2, angularly separated spectral ratios and nonzero real q,kappa with opposite inside/outside unit-circle magnitudes give absolute convergence of the nonstationary f\_N,infinity Ruijsenaars series near its affine ratio origin\end{center}
\noindent\textbf{Stable ID:} 1622\par
\noindent\textbf{Authors:} Edwin Langmann; Masatoshi Noumi; Junichi Shiraishi\par
\noindent\textbf{Original work:} Basic Properties of Non-Stationary Ruijsenaars Functions\par
\noindent\textbf{Version:} Official SIGMA 16 (2020) article 105, DOI 10.3842/SIGMA.2020.105; journal PDF and native TeX acquired 2026-09-30, exact bytes pinned by SHA256\par
\noindent\textbf{Selected task scope:} Exactly Theorem3.7 restricted to N>=2 and the f\_N,infinity series of Definitions3.1 and3.3, with s\_i nonzero, t nonzero fixed, some sigma>0 satisfying |sin(arg(s\_i/s\_j))|>sigma for all i<j, and nonzero real q,kappa such that either |q|<1<|kappa| or |kappa|<1<|q|. There exists rho>0 with absolute convergence when |p|<rho\textasciicircum{}N and all |x\_(i+1)/x\_i|<rho and |p*x\_1/x\_N|<rho; x\_i are nonzero. The series is the explicitly printed periodic matrix/partition expansion; the original function conversion and asserted eigenfunction properties are not separate selected statements.\par
\noindent\textbf{Difficulty:} H2: The key obligation is a uniform exponential bound on coefficients indexed by N partitions, despite infinitely many periodic spectral factors and potentially small denominators. The author separates angular and same-period factors, controls their integer exponents by partition inequalities and telescopes the complete product into a bound in total partition size. This provides an actual analytic convergence domain for a formal nonstationary special function and exceeds checking a simple power series by the ratio test.\par
\noindent\textbf{Editorial boundary note (not author text):} The full original convergence Theorem 3.7 and complete author proof are retained, with the selected N>=2 branch recorded on the cover. Definitions of the explicitly selected periodic f\_N,infinity series, the matrix/partition reindexing and its Appendix A.1 details, all Appendix B coefficient-estimate cases and the author-supplied Euler partition identity are included. Conversion to the separate original function and conjectural eigenfunction/duality claims are not counted.\par
\noindent\textbf{Source:} \url{https://sigma-journal.com/2020/105/}\quad DOI: \url{https://doi.org/10.3842/SIGMA.2020.105}\par
\noindent\textbf{License and attribution:} Copyright retained by the named authors. Published by SIGMA under CC BY 4.0: \url{https://creativecommons.org/licenses/by/4.0/}. Source: \url{https://sigma-journal.com/about.html}. Adaptation consists of standalone excerpt formatting, cover and numbering restoration; original author language and mathematical text are retained.\par
\noindent\textbf{Editorial symbol notice (not author text):} Editorial source notice, separate from the author: AppendixB.2.3 first product (PDF p.20/native line1041) prints exponent n-theta-m+1 in its denominator, while LemmaB.1(c)/(B.7) and the immediately following product both use the unshifted exponent n-theta-m. Those explicit adjacent expressions fix the single stray +1. Original source pages retain it; no replacement estimate or independently supplied proof is inserted.\par
\medskip\hrule\medskip\noindent\textbf{Author text begins below.}\par
\setcounter{section}{1}
\setcounter{subsection}{0}
\setcounter{subsubsection}{0}
\setcounter{Theorem}{0}
\setcounter{equation}{0}
% Original source lines 98--114
\noindent {\bf Notation:} Throughout the paper, the symbols $q$, $t$, $p$, $\kappa$ (complex parameters) and~$N$ (variable number) have special significance. We~use the following standard notation,
\begin{gather*}
(z;q)_\infty \equiv \prod_{n=0}^\infty\big(1-zq^n\big)\qquad (|q|<1),
\\
(z;q)_k\equiv \frac{(z;q)_\infty}{(q^k z;q)_\infty} \qquad (k\in\Z),
\\
(z;q,p)_\infty\equiv \prod_{n,m=0}^\infty \big(1-q^np^m z\big)\qquad (|q|<1,\ |p|<1),
\\
\theta(z;p) \equiv (z;p)_\infty(p/z;p)_\infty
\end{gather*}
for~$z\in\C$. Moreover, $T_{q,z}=q^{z\partial_z}$, i.e.,
\begin{gather*}
(T_{q,z}f)(z)=f(qz)
\end{gather*}
for functions $f(z)$ of~$z\in\C$.
For $z\in\C$, $\Re(z)$ and $\Im(z)$ are the real- and imaginary parts of~$z$, and $\sin\arg(z) = \Im(z)/|z|$. For $x=(x_1,\dots,x_N)$ and $\lambda=(\lambda_1,\dots,\lambda_N)$, $x^\lambda$ is short for~$x_1^{\lambda_1}\cdots x_N^{\lambda_N}$, $x^{-1}$ is short for~$\big(x_1^{-1},\dots,x_N^{-1}\big)$, and $x^{+1}=x$.
We~denote as $\C[[z_1,\dots,z_N]]$ the space of~all formal power series $f(z)=\sum_{\mu\in\Z_{\geq 0}^N}c_\mu z_1^{\mu_1}\cdots z_N^{\mu}$ in formal variables $z=(z_1,\dots, z_N)$ with complex coefficients~$c_\mu$.

\setcounter{section}{2}
\setcounter{subsection}{2}
\setcounter{subsubsection}{0}
\setcounter{Theorem}{3}
\setcounter{equation}{12}
% Original source lines 238--242
\section{Results on the non-stationary Ruijsenaars function}\label{sec:results}

We~give alternative series representations of~the non-stationary Ruijsenaars functions (Section~\ref{sec:alternative}) and prove convergence of~these series in a~suitable domain (Section~\ref{sec:convergence}).
\subsection{Alternative series representations}
\label{sec:alternative}

\setcounter{section}{3}
\setcounter{subsection}{1}
\setcounter{subsubsection}{0}
\setcounter{Theorem}{0}
\setcounter{equation}{0}
% Original source lines 248--277
\begin{Definition}
\label{def2}
For $N\in\Z_{\geq 1}$, two parameters $q$, $t$, and two sets of~infinitely many variables $\bx=(x_1,x_2,\dots)$ and $\bs=(s_1,s_2,\dots)$, let the following define a~formal power series in the infinitely many variables
$(x_2/x_1,x_3/x_2,x_4/x_3,\dots)$,
\begin{gather}
\label{fNinfty}
 f_{N,\infty}(\bx|\bs|q,t) \equiv \sum_{\theta\in\hat{\mathsf{M}}_N} c_{N,\infty}(\theta|\bs|q,t) \prod_{i=1}^N \prod_{k>i} (x_k/x_i)^{\theta_{ik}}
\end{gather}
with $\hat{\mathsf{M}}_N$ the set of~infinite, $N$-periodic, strictly upper triangular matrices with nonnegative integer entries which are non-zero only in a~finite strip away from the diagonal:
\begin{gather}
\label{hMMN}
\hat{\mathsf{M}}_N\equiv \big\{ \theta=(\theta_{ik})_{i,k=1}^\infty\,|\, \theta_{ik}=\theta_{i+N,k+N}\in\Z_{\geq 0}\ (i,k\geq 1), \ \theta_{ik}=0\ (k\leq i,\ k\gg i)\big\} ,
\end{gather}
and
\begin{gather}
c_{N,\infty}(\theta|\bs|q,t)
\equiv
\prod_{i=1}^N \prod_{i< j\leq k< \infty}
\frac{\big(q^{\sum_{a> k}(\theta_{ia} - \theta_{ja})}ts_{j}/s_i;q\big)_{\theta_{ik}}}{\big(q^{\sum_{a> k}(\theta_{ia} - \theta_{ja})}qs_{j}/s_i;q\big)_{\theta_{ik}}}\nonumber
\\ \hphantom{c_{N,\infty}(\theta|\bs|q,t)\equiv }
{}\times\prod_{i=1}^N \prod_{i\leq j<k<\infty}
\frac{\big(q^{-\theta_{jk}-\sum_{a>k}(\theta_{ja}-\theta_{ia})}qs_{j}/ts_i ;q\big)_{\theta_{ik}}}{\big(q^{-\theta_{jk}-\sum_{a>k}(\theta_{jb}-\theta_{ia})}s_{j}/s_i ;q\big)_{\theta_{ik}}} .
\label{cNinfty}
\end{gather}

\end{Definition}

Note that the product in~\eqref{cNinfty} always contains only a~finite number of~factors different from $1$. Moreover, by the condition $\theta_{ik}=\theta_{i+N,k+N}$,
a matrix $\theta\in\hat{\mathsf{M}}_N$ is fully determined by the matrix elements $\theta_{ik}$ for~$1\leq i\leq N$ and $1\leq k<\infty$.
Furthermore, matrices in $\mathsf{M}_N$ can be naturally identified with matrices $\theta$ in $\hat{\mathsf{M}}_N$ by setting $\theta_{ik}=0$ if $i>N$, or $k>N$, or both.

\setcounter{section}{3}
\setcounter{subsection}{1}
\setcounter{subsubsection}{0}
\setcounter{Theorem}{2}
\setcounter{equation}{4}
% Original source lines 289--292
\begin{gather}
\label{period1}
x_{i+N}=px_i,\qquad s_{i+N}=\kappa s_i \qquad (i\geq 1).
\end{gather}

\setcounter{section}{3}
\setcounter{subsection}{1}
\setcounter{subsubsection}{0}
\setcounter{Theorem}{2}
\setcounter{equation}{5}
% Original source lines 298--322
In the following, it~is sometimes convenient to use a~notation for the functions $f_{N,\infty}$ that emphasizes that the arguments $\bx$ and $\bs$ are fixed by $x$, $s$, $p$ and $\kappa$:

\begin{Definition}
\label{def:fNinfty}
We~write
%\begin{subequations}
%\label{fNinfty22}
\begin{gather*}
%\label{fNinftydef}
f_{N,\infty}(x,p|s,\kappa|q,t) \equiv f_{N,\infty}(\bx|\bs|q,t)
\end{gather*}
if $\bx\,{=}\,(x_1,x_2,\dots)$ and $\bs\,{=}\,(s_1,s_2,\dots)$ on the right-hand side are determined by $x\,{=}\,(x_1,\dots,x_N)$, $p$, $s=(s_1,\dots,s_N)$, and $\kappa$ as in~\eqref{period1}.
Thus
\begin{gather*}
f_{N,\infty}(x,p|s,\kappa|q,t) = \sum_{\theta\in\hat{\mathsf{M}}_N} c_{N,\infty}(\theta|s,\kappa|q,t)e_{N,\infty}(\theta|x,p)
\end{gather*}
with
\begin{gather}\label{cNeNdef}
c_{N,\infty}(\theta|s,\kappa|q,t) \equiv c_{N,\infty}(\theta|\bs|q,t),
\qquad
e_{N,\infty}(\theta|x,p)\equiv \prod_{i=1}^N \prod_{k=i+1}^\infty (x_k/x_i)^{\theta_{ik}}
\end{gather}
and the identifications in~\eqref{period1} on the right-hand side in~\eqref{cNeNdef}.
%\end{subequations}
\end{Definition}

\setcounter{section}{3}
\setcounter{subsection}{1}
\setcounter{subsubsection}{0}
\setcounter{Theorem}{4}
\setcounter{equation}{7}
% Original source lines 343--365
\begin{Lemma}
\label{lem1}
The formal power series in~\eqref{fNinfty}--\eqref{cNinfty} can be written as
%\begin{subequations}
%\label{fNinfty111}
\begin{gather}
\label{fNinfty1}
 f_N(\bx|\bs|q,t) = \sum_{\vlam\in\mathsf{P}^N} C_{N,\infty}(\vlam|\bs|q,t) \prod_{i=1}^N\prod_{k\geq 1}(x_{i+k}/x_{i+k-1})^{\lambda^{(i)}_k},
\end{gather}
with $\mathsf{P}^N$ the set of~all $N$-partitions $\vlam=\big(\lambda^{(1)},\lambda^{(2)},\dots,\lambda^{(N)}\big)$, $\lambda^{(i)}$ a~partition of~arbitrary length for~$i=1,\dots,N$, and
\begin{gather}
C_{N,\infty}(\vlam|\bs|q,t) =
\prod_{i=1}^N \prod_{i< j\leq k< \infty}
\frac{\big(q^{\lambda^{(i)}_{k-i+1} - \lambda^{(j)}_{k-j+1}}ts_{j}/s_i;q\big)_{\lambda^{(i)}_{k-i}-\lambda^{(i)}_{k-i+1}}}{\big(q^{\lambda^{(i)}_{k-i+1} - \lambda^{(j)}_{k-j+1}}qs_{j}/s_i;q\big)_{\lambda^{(i)}_{k-i}-\lambda^{(i)}_{k-i+1}}} \nonumber
\\ \hphantom{C_{N,\infty}(\vlam|\bs|q,t) = }
{} \times \prod_{i=1}^N \prod_{i\leq j<k<\infty}
\frac{\big(q^{-\lambda^{(j)}_{k-j}+\lambda^{(i)}_{k-i+1}}qs_{j}/ts_i ;q\big)_{\lambda^{(i)}_{k-i}-\lambda^{(i)}_{k-i+1}}}{\big(q^{-\lambda^{(j)}_{k-j}+\lambda^{(i)}_{k-i+1}}s_{j}/s_i ;q\big)_{\lambda^{(i)}_{k-i}-\lambda^{(i)}_{k-i+1}}}\label{CNinfty}
\end{gather}
setting $\lambda^{(i+N)}_j \equiv \lambda^{(i)}_j$.
\end{Lemma}
\begin{proof}
Straightforward computations, using that $\theta_{ik} = \lambda^{(i)}_{k-i}-\lambda^{(i)}_{k-i+1}$ defines a~one-to-one correspondence between multi-partitions $\vlam=\big(\lambda^{(1)},\dots,\lambda^{(N)}\big)$ in $\mathsf{P}^N$ and matrices $\theta=(\theta_{ik})_{i,k=1}^\infty$ in $\hat{\mathsf{M}}_N$ (the interested reader can find the details in Appendix~\ref{app:prooflem1}).
\end{proof}

\setcounter{section}{3}
\setcounter{subsection}{1}
\setcounter{subsubsection}{0}
\setcounter{Theorem}{6}
\setcounter{equation}{10}
% Original source lines 397--504
\subsection{Convergence}\label{sec:convergence}

We~prove that the non-stationary Ruijsenaars functions $f_N\!(x,p|s,\kappa|q,t)$ in Definitions~\ref{def2} and~\ref{def:fNinfty} are absolutely convergent in a~certain domain of~variables and parameters.

\begin{Theorem}\label{thm:converegence}
For fixed $N\in\Z_{\geq 1}$, assume that the variables $s=(s_1,\dots,s_N)\in\C^N$ and the parameters $q$ and $\kappa$ satisfy the following conditions,
\begin{itemize}\itemsep=0pt
\item[$(i)$] for some $\sigma>0$,
\begin{gather*}
{}|\sin\arg(s_i/s_j)|>\sigma \qquad (1\leq i<j\leq N),
\end{gather*}
\item[$(ii)$] $q$ and $\kappa$ both are real, and either $|q|<1$ and $|\kappa|>1$, or $|q|>1$ and $|\kappa|<1$.
\end{itemize}
Then, there exists a~constant $\rho>0$ such that the formal power series
\begin{gather*}
f_{N,\infty}(x,p;s,\kappa|q,t)\in \C[[x_2/x_1,\dots,x_N/x_{N-1},px_1/x_N]]
\end{gather*}
in Definitions~{\rm \ref{def2}} and~{\rm \ref{def:fNinfty}} is absolutely convergent in the domain
\begin{gather}\label{xcond}
|p|<\rho^N,\qquad |x_2/x_1|<\rho,\qquad \dots,\qquad |x_N/x_{N-1}|<\rho,\qquad |px_1/x_N|<\rho.
\end{gather}
\end{Theorem}
\begin{Remark}In our proof, we actually show convergence for any $\rho<1/C_1C_2$ where
\begin{gather}
C_1= 1+|1-t/q|\max\left(\frac1{\sigma}, \frac{|\kappa|}{|1-|\kappa||}\right)\!,\nonumber\\
C_2= 1+|1-q/t|\max\left(\frac1{\sigma},\frac{1}{|1-|q||}\right)\!.\label{C1C2}
\end{gather}
\end{Remark}
\begin{Remark}
We~believe that it~is possible to refine this convergence result. In particular, we~believe that there are regions of~convergence where $s_i/s_j$, $1\leq i<j\leq N$, are real and $q$ and~$\kappa$ have non-trivial imaginary parts.
\end{Remark}

\begin{proof}[Proof of~Theorem~\ref{thm:converegence}]
Our strategy of~proof is to show that our assumptions imply simple upper bounds on the terms appearing in the series in~\eqref{fNinfty1}--\eqref{CNinfty}:
\begin{gather}
\label{estim}
\Bigg|\prod_{i=1}^N\prod_{k\geq 1}(x_{i+k}/x_{i+k-1})^{\lambda^{(i)}_k}\Bigg|\leq \rho^{|\vlam|},\qquad \left|C_{N,\infty}(\vlam|\bs|q,t)\right|\leq C_1^{|\vlam|}C_2^{|\vlam|}
\end{gather}
with $|\vlam|\equiv \sum_{i=1}^N \sum_{k\geq 1}\lambda^{(i)}_k$ and $\alpha = \rho C_1C_2<1$.
With that, absolute convergence follows from the comparison test: the series in~\eqref{fNinfty1}--\eqref{CNinfty} is of~the form $\sum_{\vlam\in\mathsf{P}^N} a_{\vlam}$ with $|a_{\vlam}|\leq \alpha^{|\vlam|}$ for all $\vlam\in\mathsf{P}^N$,
and the series $\sum_{\vlam\in\mathsf{P}^N} \alpha^{|\vlam|}$ converges absolutely for~$|\alpha|<1$.

The first estimate in~\eqref{estim} is a~simple consequence of~the conditions in~\eqref{xcond}: since \allowbreak \mbox{$x_{i+N}=px_i$} for all $i\geq 1$, these conditions are equivalent to
\begin{gather*}
|x_{i+1}/x_i| < \rho\qquad (i\geq 1) ,
\end{gather*}
which clearly implies the result.

The proof of~the second estimate in~\eqref{estim} is more involved and, for this reason, we supplement our somewhat descriptive arguments in the main text below by a~detailed argument in~Appendix~\ref{app:estim}.

We~observe that $C_{N,\infty}(\vlam|\bs|q,t)$ in~\eqref{CNinfty} is a~product of~fractions $\big(1-q^{l}au\big)/\big(1-q^l u\big)$ with $a=t/q$ in the first group of~products and $a=q/t$ in the second group,
 $l\in\Z$, and $u=s_j/s_i$ for~$i=1,\dots,N$ and $j\geq i$; moreover, $s_{j+\ell N}=\kappa^\ell s_j$ for~$\ell\in\Z_{\geq 1}$.
Such a~fraction can be estimated in a~simple way:
\begin{gather*}
\left|\frac{1-q^{l}au}{1-q^l u}\right| = \left|1+(1-a)\frac{q^l u}{1-q^l u} \right|\leq 1+|1-a|\left|\frac{q^l u}{1-q^l u}\right|\!.
\end{gather*}
If $j-i$ is {\em not} an integer multiple of~$N$, we can estimate this further using
\begin{gather}
\label{zest}
\left| \frac{z}{1-z}\right| \leq \frac1{|\sin\arg(z)|} \qquad (z\in\C\setminus\{\R\})
\end{gather}
(to see that the latter inequality holds, write $z=|z|\ee^{\ii\varphi}$ and note that~\eqref{zest} is equivalent to
\begin{gather*}
|z|^2\sin^2\varphi \leq 1+|z|^2-2|z|\cos\varphi\Leftrightarrow 0\leq (1-|z|\cos\varphi)^2,
\end{gather*}
which is obvious). Since we assume that $q$ and $\kappa$ both are real,
\begin{gather*}
\big|\sin\arg\big(q^l \kappa^{\ell }s_j/s_i\big)\big| = |\sin\arg(s_j/s_i)| \geq \sigma>0\qquad (j-i\neq N\Z_{\geq 0})
\end{gather*}
for all integers $l,\ell$, we get a~simple universal bound for these fractions:
\begin{gather*}
\left|\frac{1-q^l a s_j/s_i}{1-q^l s_j/s_i}\right| \leq 1+|1-a|\frac{1}{\sigma} \qquad (j-i \notin N\Z_{\geq 0})
\end{gather*}
for all integers $l$. However, this bound does not work for~$j=i+\ell N$ with $\ell\in\Z_{\geq 0}$ since, in these cases, $q^l u=q^l s_j/s_i=q^l \kappa^\ell$ is real.
However, one can check that, in all these latter cases, either $l\leq 0$ and $\ell>0$, or $l<0$ and $\ell\geq 0$, and thus, by our assumptions, $z\equiv q^lu=q^l \kappa^\ell$ {\em always}
satisfies either $|z|\geq \min\big(|q|^{-1},|\kappa|\big)>1$ (if~$|q|<1$ and $|\kappa|>1$) or $|z|\leq \max\big(|q|,|\kappa|^{-1}\big)<1$ (if~$|q|>1$ and $|\kappa|<1$); we therefore can use the inequality
\begin{gather*}
%\label{zest1}
\left| \frac{z}{1-z}\right| \leq \frac{|z|}{|1-|z||} \qquad (|z|\neq 1)
\end{gather*}
to get simple universal bounds for the cases $j=i+\ell N$ with $\ell\in\Z_{\geq 0}$ as well (we spell our the details of~this argument in Appendix~\ref{app:estb}).
We~thus get estimates
\begin{gather*}
\left|\frac{1-q^{l}au}{1-q^l u}\right| \leq C_{1,2}
\end{gather*}
with different upper bounds, $C_1$ and $C_2$, for {\em all} fractions in the first and second groups of~pro\-ducts on the right-hand side in~\eqref{CNinfty}, respectively.
The arguments above allow to compute the constants $C_1$ and $C_2$ and give the results in~\eqref{C1C2}; the interested reader can find the details of~this computation in Appendix~\ref{app:estim}.

Inserting these bounds into~\eqref{CNinfty} we obtain
\begin{gather}
|C_{N,\infty}(\vlam|\bs|q,t)| \leq
\prod_{i=1}^N \Bigg( \prod_{i< j\leq k< \infty} C_1^{\lambda^{(i)}_{k-i}-\lambda^{(i)}_{k-i+1}}\Bigg)\Bigg( \prod_{i\leq j< k< \infty} C_2^{\lambda^{(i)}_{k-i}-\lambda^{(i)}_{k-i+1}}\Bigg) \nonumber
\\ \hphantom{|C_{N,\infty}(\vlam|\bs|q,t)|}
{}= \prod_{i=1}^N \Bigg( \prod_{i< j< \infty} C_1^{\lambda^{(i)}_{j-i}}\Bigg)\Bigg( \prod_{i\leq j< \infty} C_2^{\lambda^{(i)}_{j+1-i}}\Bigg)\nonumber
\\ \hphantom{|C_{N,\infty}(\vlam|\bs|q,t)|}
{}= \prod_{i=1}^N \Bigg( \prod_{k\geq 1} C_1^{\lambda^{(i)}_{k}}\Bigg) \Bigg( \prod_{k\geq 1} C_2^{\lambda^{(i)}_{k}}\Bigg) = C_1^{|\vlam|}C_2^{|\vlam|},
\label{Cbound}
\end{gather}
computing telescoping products in the second step and using $\sum_{i=1}^N\sum_{k\geq 1}\lambda^{(i)}_k=|\vlam|$ in the last step.
This proves the second estimate in~\eqref{estim}.

To conclude, we prove that the series $\sum_{\vlam\in\mathsf{P}^N} \alpha^{|\vlam|}$ for~$|\alpha|<1$ is absolutely convergent by the following computation,
\begin{gather*}
\sum_{\vlam\in\mathsf{P}^N} \alpha^{|\vlam|} = \sum_{\lambda^{(1)},\dots,\lambda^{(N)}\in\mathsf{P}} \prod_{i=1}^N \alpha^{|\lambda^{(i)}|} =
 \prod_{i=1}^N \sum_{\lambda^{(i)}\in\mathsf{P}} \alpha^{|\lambda^{(i)}|} =\left( \sum_{\lambda\in\mathsf{P}} \alpha^{|\lambda|}\right)^N = \frac1{(\alpha;\alpha)_\infty^N},
\end{gather*}
using the definition $|\lambda|\equiv \sum_{k\geq 1}\lambda_k$ for partitions $\lambda$; for clarity, and for the convenience of~the reader, we give in Appendix~\ref{app:Id} the well-known identity used in the last step, together with its elementary proof making absolute convergence manifest; see~\eqref{Euler}--\eqref{Euler1}.
\end{proof}

\setcounter{section}{5}
\setcounter{subsection}{0}
\setcounter{subsubsection}{0}
\setcounter{Theorem}{0}
\setcounter{equation}{0}
% Original source lines 741--744
\appendix

\section{Alternative series representation}
\label{app:proof}

\setcounter{section}{1}
\setcounter{subsection}{0}
\setcounter{subsubsection}{0}
\setcounter{Theorem}{0}
\setcounter{equation}{0}
% Original source lines 748--770
\subsection[Details on Lemma~3.5]{Details on Lemma~\ref{lem1}}\label{app:prooflem1}
One can check that the following two formulas provide a~correspondence between multi-partitions $\vlam=\big(\lambda^{(1)},\dots,\lambda^{(N)}\big)$ in $\mathsf{P}^N$ and matrices $\theta=(\theta_{ik})_{i,k=1}^\infty$ in $\hat{\mathsf{M}}_N$ that is one-to-one:
\begin{gather*}
%\label{tetfromlam}
\theta_{ik} = \lambda^{(i)}_{k-i}-\lambda^{(i)}_{k-i+1}
\end{gather*}
and
\begin{gather}\label{lamfromtet}
\lambda^{(i)}_{k-i} = \sum_{a\geq k}\theta_{ia}
\end{gather}
setting $\lambda^{(i+N)}=\lambda^{(i)}$ for all $i\geq 1$. With this identification, $C_{N,\infty}(\vlam|\bs|q,t)$ in~\eqref{CNinfty} is clearly equal to $c_{N,\infty}(\theta|\bs|q,t)$ in~\eqref{cNinfty}, and
\begin{gather*}
\begin{split}
&\prod_{i=1}^N \prod_{k=i+1}^\infty (x_k/x_i)^{\theta_{ik}} = \prod_{i=1}^N \prod_{k=i+1}^\infty [(x_k/x_{k-1})(x_{k-1}/x_{k-2})\cdots (x_{i+1}/x_i)]^{\theta_{ik}}
\\ &\hphantom{\prod_{i=1}^N \prod_{k=i+1}^\infty (x_k/x_i)^{\theta_{ik}}}
{}= \prod_{i=1}^N \prod_{i<j\leq k<\infty} (x_j/x_{j-1})^{ \lambda^{(i)}_{k-i}-\lambda^{(i)}_{k-i+1}}
 = \prod_{i=1}^N \prod_{j>i} (x_j/x_{j-1})^{\lambda^{(i)}_{j-i}}
\\ &\hphantom{\prod_{i=1}^N \prod_{k=i+1}^\infty (x_k/x_i)^{\theta_{ik}}}
{}= \prod_{i=1}^N \prod_{k\geq 1}(x_{i+k}/x_{i+k-1})^{\lambda^{(i)}_k},
\end{split}
\end{gather*}
inserting a~telescoping product in the second step, using~\eqref{lamfromtet} in the third, and computing a~telescoping product in the fourth. This proves the result.


\setcounter{section}{1}
\setcounter{subsection}{2}
\setcounter{subsubsection}{0}
\setcounter{Theorem}{0}
\setcounter{equation}{2}
% Original source lines 894--1052
\section{Estimates}\label{app:estim}
{\sloppy
We~give a~complementary proof of~the second estimate in~\eqref{estim}, to compute the upper bounds~$C_{1,2}$ in Theorem~\ref{thm:converegence}.

}

\subsection[Complementary proof of~the second estimate in (3.13)]{Complementary proof of~the second estimate in~(\ref{estim})}
\label{app:proofestim}
We~prove that, under the assumptions in Theorem~\ref{thm:converegence}, the following estimates hold true for the fractions appearing in the formula~\eqref{CNinfty} for~$C_{N,\infty}(\vlam|\bs|q,t)$,
\begin{gather}
\label{frac1}
\left|\frac{\big(q^{\lambda^{(i)}_{k-i+1} - \lambda^{(j)}_{k-j+1}}ts_{j}/s_i;q\big)_{\lambda^{(i)}_{k-i}-\lambda^{(i)}_{k-i+1}}}{\big(q^{\lambda^{(i)}_{k-i+1} - \lambda^{(j)}_{k-j+1}}qs_{j}/s_i;q\big)_{\lambda^{(i)}_{k-i}-\lambda^{(i)}_{k-i+1}}}\right|\leq C_1^{\lambda^{(i)}_{k-i}-\lambda^{(i)}_{k-i+1}} \ \ (1\leq i\leq N,\ i<j\leq k<\infty) ,\!\!\!
\\
\label{frac2}
\left| \frac{\big(q^{-\lambda^{(j)}_{k-j}+\lambda^{(i)}_{k-i+1}}qs_{j}/ts_i ;q\big)_{\lambda^{(i)}_{k-i}-\lambda^{(i)}_{k-i+1}}}{\big(q^{-\lambda^{(j)}_{k-j}+\lambda^{(i)}_{k-i+1}}s_{j}/s_i ;q\big)_{\lambda^{(i)}_{k-i}-\lambda^{(i)}_{k-i+1}}} \right| \leq C_2^{\lambda^{(i)}_{k-i}-\lambda^{(i)}_{k-i+1}} \ \
(1\leq i\leq N,\ i\leq j< k<\infty)\!\!
\end{gather}
for all $N$-partitions $\vlam=\big(\lambda^{(1)},\dots,\lambda^{(N)}\big)\in\mathsf{P}^N$, with $C_1$ and $C_2$ in~\eqref{C1C2}.
This and~\eqref{CNinfty} imply the estimate in~\eqref{Cbound} which, by the computation in~\eqref{Cbound}, is equivalent to the second estimate in~\eqref{estim}.

We~observe all estimates in~\eqref{frac1}--\eqref{frac2} are of~the form
\begin{gather*}
\left|\frac{\big(q^{l}as_i/s_j;q\big)_\theta}{\big(q^l s_i/s_j;q\big)_\theta}\right|\leq C^\theta,
\end{gather*}
where
\begin{gather}
\label{first}
l= \lambda^{(i)}_{k-i+1} - \lambda^{(j)}_{k-j+1}+1,\qquad
a=t/q,\qquad
\theta= \lambda^{(i)}_{k-i}-\lambda^{(i)}_{k-i+1},\qquad
C=C_1
\end{gather}
in~\eqref{frac1} and
\begin{gather}
\label{second}
l=-\lambda^{(j)}_{k-j}+\lambda^{(i)}_{k-i+1},\qquad
a=q/t,\qquad
\theta= \lambda^{(i)}_{k-i}-\lambda^{(i)}_{k-i+1},\qquad
C=C_2
\end{gather}
in~\eqref{frac2}. We~prove~\eqref{frac1}--\eqref{frac2} using three different kinds of~estimates:

\begin{Lemma}\label{lem:est}
Let $\theta\in\Z_{\geq 0}$, $a\in\C$, $q, \kappa\in\R$ with either $|q|<1$ and $|\kappa|>1$ or $|q|>1$ and $|\kappa|<1$. Then the following estimates hold true,
\begin{itemize}\itemsep=0pt
\item[$(a)$] for all $l\in\Z$ and $u\in\C\setminus\{\R\}$:
\begin{gather}\label{esta}
\left| \frac{\big(q^{l}au;q\big)_\theta}{\big(q^l u;q\big)_\theta}\right|\leq \left( 1 + \frac{|1-a|}{|\sin\arg(u)|}\right)^\theta\!,
\end{gather}
\item[$(b)$] for all $m\in\Z_{\geq 0}$, $\ell\in\Z_{\geq 1}$:
\begin{gather}\label{estb}
\left|\frac{\big( q^{-\theta-m+1} a \kappa^\ell \big)_\theta}{\big(q^{-\theta-m+1} \kappa^\ell \big)_\theta} \right| \leq \left(1+|1-a|\frac{|\kappa|}{|1-|\kappa||} \right)^\theta\!,
\end{gather}
\item[$(c)$] for all $m\in\Z_{\geq 0}$, $\ell\in\Z_{\geq 0}$:
\begin{gather}\label{estc}
\left|\frac{\big( q^{-\theta-m} a \kappa^\ell \big)_\theta}{\big(q^{-\theta-m} \kappa^\ell \big)_\theta} \right| \leq \left(1+|1-a|\frac{1}{|1-|q||} \right)^\theta\!.
\end{gather}
\end{itemize}
\end{Lemma}
(The proof is given in Appendix~\ref{app:prooflem:est}.)

{\bf Case A:} For $j-i\notin N\Z_{\geq 0}$, we can use the estimate in~\eqref{esta}:
Since $s_{j+N}=\kappa s_j$ and~$\kappa$ is real, we have $|\sin\arg(s_j/s_i)|=|\sin\arg(s_{j+N}/s_i)|=|\sin\arg(s_i/s_j)|$ for all $j\geq i$;
since $|\sin\arg(s_j/s_i)|\geq \sigma$ for all $1\leq i<j\leq N$ by assumption, $|\sin\arg(s_j/s_i)|\geq \sigma$ for all $1\leq i\leq N$ and $j\geq i$ such that $j-i\neq N\Z_{\geq 0}$, and we get
\begin{gather*}
\left|\frac{(q^{l}as_i/s_j;q)_\theta}{(q^l s_i/s_j;q)_\theta}\right| \leq \left(1+\frac{|1-a|}{\sigma} \right)^\theta\qquad
(j-i\notin N\Z_{\geq 0})
\end{gather*}
for all cases in~\eqref{first}--\eqref{second}. This proves that the estimates in~\eqref{frac1}--\eqref{frac2} for all
\begin{gather}
\label{C1C2A}
C_1\geq 1+\frac{|1-t/q|}{\sigma},\qquad
C_2\geq 1+\frac{|1-q/t|}{\sigma}
\end{gather}
and for all cases $j-i\notin N\Z_{\geq 0}$.

We~consider the remaining cases for~\eqref{frac1} and~\eqref{frac2} below in Cases B and C, respectively.

{\bf Case B:} For $j-i\in N\Z_{\geq 1}$, we have $s_j/s_i = s_{i+\ell N}/s_i=\kappa^\ell$ for some $\ell\in\Z_{\geq 1}$, and we can use the estimate in~\eqref{estb}:
\begin{gather*}
\left|\frac{\big(q^{-\theta-m+1}as_i/s_j;q\big)_\theta}{\big(q^{-\theta-m+1} s_i/s_j;q\big)_\theta}\right|\leq \left(1+|1-a|\frac{|\kappa|}{||\kappa|-1|} \right)^\theta\qquad
(j-i\in N\Z_{\geq 1},\ m\in\Z_{\geq 0}).
\end{gather*}
We~check that all cases in~\eqref{frac1} for~$j-i\in N\Z_{\geq 1}$ are covered by this: all $l$ in~\eqref{first} for~$j=i+\ell N$ can be written as (recall that $\lambda^{(i+\ell N)}_k=\lambda^{(i)}_k$)
\begin{gather*}
l = -\big[\lambda^{(i)}_{k-i}- \lambda^{(i)}_{k-i+1}\big] - \big[\lambda^{(i)}_{k-i-\ell N+1}-\lambda^{(i)}_{k-i}\big]+1 = -\theta-m+1
\end{gather*}
with $m=\lambda^{(i)}_{k-i-\ell N+1}-\lambda^{(i)}_{k-i}\geq 0$ since $\ell\geq 1$ and $\lambda^{(i)}=\big(\lambda_1^{(i)},\lambda_2^{(i)},\dots\big)$ is a~partition.
This proves that~\eqref{frac1} holds true if
\begin{gather}
\label{C1C2B}
C_1\geq 1+|1-t/q|\frac{|\kappa|}{|1-|\kappa||}
\end{gather}
for all cases $j-i\in N\Z_{\geq 1}$.

{\bf Case C:} For $j-i\in N\Z_{\geq 0}$, we have $s_j/s_i = s_{i+\ell N}/s_i=\kappa^\ell$ for some $\ell\in\Z_{\geq 0}$, and we can use the estimate in~\eqref{estc}:
\begin{gather*}
\left|\frac{\big(q^{-\theta-m}as_i/s_j;q\big)_\theta}{\big(q^{-\theta-m} s_i/s_j;q\big)_\theta}\right|\leq \left(1+|1-a|\frac{1}{|1-|q||} \right)^\theta\qquad
(j-i\in N\Z_{\geq 0},\ m\in\Z_{\geq 0}).
\end{gather*}
We~check that all cases in~\eqref{frac2} for~$j-i\in N\Z_{\geq 0}$ are covered by this: all $l$ in can be written as
\begin{gather*}
l=-\big[\lambda^{(i)}_{k-i}-\lambda^{(i)}_{k-i+1}\big]-\big[\lambda^{(i)}_{k-i-\ell N}-\lambda^{(i)}_{k-i}\big] = -\theta-m
\end{gather*}
with $m=\lambda^{(i)}_{k-i-\ell N}-\lambda^{(i)}_{k-i}\geq 0$.
This proves that~\eqref{frac2} holds true if
\begin{gather}\label{C1C2C}
C_2\geq 1+|1-q/t|\frac{1}{|1-|q||}
\end{gather}
for all cases $j-i\in N\Z_{\geq 0}$.

We~proved that~\eqref{frac1}--\eqref{frac2} holds true for {\em all} cases provided the conditions in~\eqref{C1C2A},~\eqref{C1C2B} and~\eqref{C1C2C} all hold true; this is the case if we choose $C_1$ and $C_2$ as in~\eqref{C1C2}.

\subsection[Proof of~Lemma~B1]{Proof of~Lemma~\ref{lem:est}}\label{app:prooflem:est}
\subsubsection[Proof of~the estimate in (B.5)]{Proof of~the estimate in~(\ref{esta})}
We~have
\begin{gather*}
\text{LHS} = \left|\prod_{n=0}^{\theta-1}\frac{1-q^{l+n}au}{1-q^{l+n}u} \right| = \prod_{n=0}^{\theta-1}\left|1+(1-a)\frac{q^{l+n}u}{1-q^{l +n}u} \right|
\\ \hphantom{\text{LHS}}
{}\leq \prod_{n=0}^{\theta-1}\left(1+|1-a|\left|\frac{q^{l+n}u}{1-q^{l+n}u}\right| \right)\leq \prod_{n=0}^{\theta-1}\left(1+\frac{|1-a|}{|\sin\arg(u)|}\right) =
\text{RHS},
\end{gather*}
using the estimate in~\eqref{zest} and $\big|\sin\arg\big(q^{l+n}u\big)\big|=|\sin\arg(u)|$ since $q$ is real.

\subsubsection[Proof of~the estimate in (B.6)]{Proof of~the estimate in~(\ref{estb})}
\label{app:estb}
We~have
\begin{gather*}
\text{LHS} = \left| \prod_{n=0}^{\theta-1} \frac{1- q^{n-\theta-m+1}a\kappa^\ell }{1-q^{n-\theta-m+1} \kappa^\ell } \right| = \prod_{n=0}^{\theta-1}\left| 1+(1-a)\frac{q^{-n-m}\kappa^\ell}{1-q^{-n-m}\kappa^\ell}\right|
\\ \hphantom{\text{LHS}}
{}\leq \prod_{n=0}^{\theta-1}\left( 1+|1-a|\left|\frac{q^{-n-m}\kappa^\ell}{1-q^{-n-m}\kappa^\ell}\right|\right)
\leq \prod_{n=0}^{\theta-1}\left( 1+|1-a|\frac{|\kappa|}{|1-|\kappa||}\right) =\text{RHS},
\end{gather*}
using
\begin{gather*}
\left|\frac{q^{-l}\kappa^\ell}{1-q^{-l}\kappa^\ell}\right|\leq \frac{|\kappa|}{|1-|\kappa||}\qquad
(l\geq 0, \ell\geq 1);
\end{gather*}
the latter follows for the case $|\kappa|<1$ and $|q|>1$ from the following inequality: $x/(1-x)<y/(1-x)$ for~$0\leq x<y<1$, and for the case $|\kappa|>1$ and $|q|<1$:
\begin{gather*}
\left|\frac{q^{-l}\kappa^\ell}{1-q^{-l}\kappa^\ell}\right| = \left|\frac{1}{1-q^{l}\kappa^{-\ell}}\right| \leq \frac1{1-|1/\kappa|} = \frac{|\kappa|}{|1-|\kappa||}
\end{gather*}
since $1/(1-x)<1/(1-y)$ for~$0\leq x<y<1$.

\subsubsection[Proof of~the estimate in (B.7)]{Proof of~the estimate in~(\ref{estc})}
We~have
\begin{gather*}
\text{LHS} = \left| \prod_{n=0}^{\theta-1} \frac{1- q^{n-\theta-m}a\kappa^\ell }{1-q^{n-\theta-m+1} \kappa^\ell } \right| = \prod_{n=1}^{\theta}\left| 1+(1-a)\frac{q^{-n-m}\kappa^\ell}{1-q^{-n-m}\kappa^\ell}\right|
\\ \hphantom{\text{LHS}}
{}\leq \prod_{n=1}^{\theta}\left( 1+|1-a|\left|\frac{q^{-n-m}\kappa^\ell}{1-q^{-n-m}\kappa^\ell}\right|\right)
\leq \prod_{n=1}^{\theta}\left( 1+|1-a|\frac{1}{|1-|q||}\right) =\text{RHS},
\end{gather*}
using
\begin{gather*}
\left|\frac{q^{-l}\kappa^\ell}{1-q^{-l}\kappa^\ell}\right|\leq \frac{|q^{-1}|}{||q^{-1}|-1|} = \frac{1}{|1-|q||} \qquad
(l\geq 1,\ \ell\geq 0),
\end{gather*}
as in the proof of~\eqref{estb}.


\setcounter{section}{3}
\setcounter{subsection}{2}
\setcounter{subsubsection}{0}
\setcounter{Theorem}{1}
\setcounter{equation}{10}
% Original source lines 1240--1262
\section{Identities}\label{app:Id}
For clarify, and for the convenience of~the reader, we state and prove two identities used in the main text.

First, the identity
\begin{gather}\label{Euler}
\sum_{\lambda\in\mathsf{P}}\alpha^{|\lambda|} = \frac1{(\alpha;\alpha)_\infty} \qquad (|\alpha|<1),
\end{gather}
which goes back to Euler, is important in our proof of~Theorem~\ref{thm:converegence}. It is proved by the following elementary computation making absolute convergence of~the series manifest,
\begin{gather}
\text{LHS} = \lim_{M\to\infty} \sum_{\lambda_1\geq \lambda_2\geq \cdots\geq \lambda_M\geq 0} \alpha^{\lambda_1+\lambda_2+\dots+\lambda_M}\nonumber
\\ \phantom{\text{LHS}}
{}= \lim_{M\to\infty} \sum_{\lambda_1=\lambda_2}^\infty \sum_{\lambda_2=\lambda_3}^\infty \cdots \sum_{\lambda_{M-1}=\lambda_M}^\infty \sum_{\lambda_M=0}^\infty \alpha^{\lambda_1+\lambda_2+\dots+\lambda_M}\nonumber
\\ \phantom{\text{LHS}}
{}= \lim_{M\to\infty} \sum_{\lambda_2=\lambda_3}^\infty \cdots \sum_{\lambda_{M-1}=\lambda_M}^\infty \sum_{\lambda_M=0}^\infty \frac{1}{1-\alpha}
\alpha^{2\lambda_2+\lambda_3 +\dots+\lambda_M}\nonumber
\\ \phantom{\text{LHS}}
{}= \lim_{M\to\infty} \sum_{\lambda_3=\lambda_4}^\infty \cdots \sum_{\lambda_{M-1}=\lambda_M}^\infty \sum_{\lambda_M=0}^\infty \frac{1}{(1-\alpha)(1-\alpha^2)}
\alpha^{3\lambda_3+\lambda_4 +\dots+\lambda_M} = \cdots \nonumber
\\ \phantom{\text{LHS}}
{}= \lim_{M\to\infty} \frac1{(1-\alpha)} \frac1{(1-\alpha^2)}\cdots \frac1{(1-\alpha^M)} = \text{RHS},
\label{Euler1}
\end{gather}
summing repeatedly the geometric series.
\begin{thebibliography}{99}
\end{thebibliography}
\end{document}
